\documentclass{ximera}

\input{../../preamble.tex}

\title{Core Analysis}

\begin{document}

\begin{abstract}
Part 1
\end{abstract}
\maketitle


The Core Functions are our basic, beginning function building blocks. Linear combinations and linear compositions create the Elementary Functions.

The Core Functions cannot be broken down any further and their analysis is generally straightforward.

As a starting point for function analysis, let's just list characteristics of our Core Functions. A full analysis would also provide reasoning.


\begin{explanation} \textbf{\textcolor{red!75!green}{Constant Functions}} 


\[ Constant(x) = C  \]



\textbf{Domain:} The natural domain of every constant function is always $(-\infty, \infty)$.

\textbf{Zeros:} If the constant function just happens to be the zero function (every value is $0$), then every number in the domain is a zero of the function.

Otherwise, constant functions do not have zeros.

\textbf{Continuity:} Constant functions are always continuous.

\textbf{End-Behavior:} The end-behavior is just the value of constant function.


\textbf{Behavior (Inc/Dec):} According to the definition of increasing and decreasing, constant functions technically both increase and they decrease. However, we don't view them in this way. We view them as remaining constant.

\textbf{Global Maximum and Minimum:} Since a constant function has only one value, that value is the global maxmum and global minimum and they both occur at every domain number.

\textbf{Local Maximum and Minimum:}  Since a constant function has only one value, that value is a local maxmum and local minimum and they both occur at every domain number.

\textbf{Range:} The range of a constant function consists of just one number, the value of the function.

\end{explanation}









\begin{explanation} \textbf{\textcolor{red!75!green}{The Identity Function}} 


\[ Id(x) = x  \]



\textbf{Domain:} The natural domain of the identity function is $(-\infty, \infty)$.

\textbf{Zeros:} $0$ is the only zero of the identity function.

\textbf{Continuity:} The identity function is a continuous function.

\textbf{End-Behavior:} 

As the domain becomes unbounded positively, the identity function also becomes unbounded positively.  Our notation for this looks like

\[
\lim\limits_{x \to \infty} Id(x) = \infty
\]

As the domain becomes unbounded negatively, the identity function also becomes unbounded negatively.  Our notation for this looks like

\[
\lim\limits_{x \to -\infty} Id(x) = -\infty
\]




\textbf{Behavior (Inc/Dec):} The identity is an increasing function.

\textbf{Global Maximum and Minimum:} The identity function has no global maximum or global minimum, which is seen in the end-behavior.




\textbf{Local Maximum and Minimum:} The identity function has no local maximums or global minimums. 

\textbf{Range:} The range of the identity function is $(-\infty, \infty)$.

\end{explanation}













\begin{explanation} \textbf{\textcolor{red!75!green}{Power Functions}} 


The power makes a big difference here.  As a beginning, we will only consider two types of powers: positive integers and negative integers,


\[ Power(x) = x^n  \]

We group fractional powers in with radical functions.

\textbf{Domain:} 

If $n$ is a whole number, then the natural domain is $(-\infty, \infty)$.

If $n$ is a negative integer, then the natural domain is $(-\infty, 0) \cup (0, \infty)$.


\textbf{Zeros:} 

If $n$ is a whole number, then $0$ is the only zero.

If $n$ is a negative integer, then the power function has no zeros.


\textbf{Continuity:} Power functions are continuous functions.

\textbf{End-Behavior:} 

If $n$ is a whole number, then there are two different cases.

\begin{itemize}
\item If $n$ is odd, then the power function is unbounded positively as the domain becomes unbounded positively or negative.
\item If $n$ is even, then the power function is unbounded positively in both directions.
\end{itemize}

If $n$ is a negative integer, then the power function tends to $0$ as the domain becomes unbounded positively or negatively.



\textbf{Behavior (Inc/Dec):} 

If $n$ is a whole number, then there are two different cases.

\begin{itemize}
\item If $n$ is odd, then the power function is an increasing function.
\item If $n$ is even, then the power function decreases on $(-\infty, 0)$ and increases on $(0, \infty)$.
\end{itemize}

If $n$ is a negative integer, then there are two different cases.

\begin{itemize}
\item If $n$ is odd, then the power function is a decreasing function.
\item If $n$ is even, then the power function increases on $(-\infty, 0)$ and decreases on $(0, \infty)$.
\end{itemize}



\textbf{Global Maximum and Minimum:} 

A power function has a global minimum only if the power is an even whole number. In that case, $0$ is the global minimum, which occurs at $0$.

\textbf{Local Maximum and Minimum:} 

The only local extrema is a global minimum, if the power is even.

\textbf{Range:} 

If $n$ is a whole number, then there are two different cases.

\begin{itemize}
\item If $n$ is odd, then the range is $(-\infty, \infty)$.
\item If $n$ is even, then the range is $[0, \infty)$.
\end{itemize}

If $n$ is a negative integer, then there are two different cases.

\begin{itemize}
\item If $n$ is odd, then the range is $(-\infty, 0) \cup (0, \infty)$.
\item If $n$ is even, then the range is $(0, \infty)$.
\end{itemize}


\end{explanation}










\begin{explanation} \textbf{\textcolor{red!75!green}{Root Functions}} 


\[ Root(x) = \sqrt[n]{x}  \]

To begin, we are only considering $n$ to be a natural number.


\textbf{Domain:} 


\begin{itemize}
\item If $n$ is odd, then the natural domain is $(-\infty, \infty)$.
\item If $n$ is even, then the natural domain is $[0. \infty)$.
\end{itemize}


\textbf{Zeros:} Core root functions only have $0$ as their only zero.

\textbf{Continuity:}  Core root functions are continuous functions.

\textbf{End-Behavior:} 


\begin{itemize}
\item If $n$ is odd, then the root function is unbounded positively as the domain becomes unbounded positively and unbounded negatively as the domain becomes unbounded negatively.

Our notation for this looks like 


\[
\lim\limits_{x \to \infty} \sqrt[n]{x} = \infty
\]


\[
\lim\limits_{x \to -\infty} \sqrt[n]{x} = -\infty
\]


\item If $n$ is even, then there is only end-behavior in the positive direction.

\[
\lim\limits_{x \to \infty} \sqrt[n]{x} = \infty
\]



\end{itemize}


\textbf{Behavior (Inc/Dec):}  Core root functions are increasing functions.

\textbf{Global Maximum and Minimum:}  

\begin{itemize}
\item If $n$ is even, then the root function has $0$ as its global minimum value and no maximum global value.
\item If $n$ is odd, then there is neither a global minimum nor global maximum.
\end{itemize} 

\textbf{Local Maximum and Minimum:} 


\begin{itemize}
\item If $n$ is even, then the root function has $0$ as its local minimum value and no local maximum values.
\item If $n$ is odd, then there are no local extrema.
\end{itemize} 


\textbf{Range:} 


\begin{itemize}
\item If $n$ is even, then the range is $[0, \infty)$.
\item If $n$ is odd, then the range is $(-\infty, \infty)$.
\end{itemize} 

\end{explanation}




\textbf{Power Functions} can have any real number as the power.


\[  pow(x) = x^r  \]

Our first investigation is with integer powers.

We expand this to rational numbers through roots and radicals.

\[ pow(x) = x^{\tfrac{n}{d}} = \sqrt[d]{x^n}  \]

We'll leave irrational powers for Calculus.

























































\begin{center}
\textbf{\textcolor{green!50!black}{ooooo-=-=-=-ooOoo-=-=-=-ooooo}} \\

more examples can be found by following this link\\ \link[More Examples of Visual Behavior]{https://ximera.osu.edu/csccmathematics/precalculus1/precalculus1/visualBehavior/examples/exampleList}

\end{center}





\end{document}
